\section*{Appendix B: Homological Algebra Essentials}
\addcontentsline{toc}{section}{Appendix B: Homological Algebra Essentials}

\section{Chain Complexes}

Let $\mathbf{Vect}$ denote the category of real vector spaces. A chain complex $(C_\bullet,d)$ is a sequence of vector spaces
[
\cdots \longrightarrow C_{n+1}\xrightarrow{d_{n+1}} C_n\xrightarrow{d_n} C_{n-1}\longrightarrow\cdots
]
such that $d_n\circ d_{n+1}=0$ for all $n\in\mathbb{Z}$.  One defines the $n^{\mathrm{th}}$ homology as
[
H_n(C_\bullet)=\frac{\ker(d_n)}{\operatorname{im}(d_{n+1})}.
]

\begin{remark}
Tangent complexes of derived stacks are examples of chain complexes whose homology measures infinitesimal automorphisms and obstructions.
\end{remark}

\section{Double Complexes and Totalization}

A double complex $C_{\bullet,\bullet}$ is a bi-graded collection ${C_{p,q}}$ with horizontal and vertical differentials. The total complex $\operatorname{Tot}(C)$ is defined by
[
\operatorname{Tot}(C)*n = \bigoplus*{p+q=n} C_{p,q}.
]
Totalization is crucial in defining the cotangent complex and for bookkeeping in BV--BRST constructions.

\section{Tensor Products and Hom Functors}

For chain complexes $C_\bullet$ and $D_\bullet$, the tensor product is the complex
[
(C\otimes D)*n = \bigoplus*{p+q=n} C_p\otimes D_q,
]
with differential determined by the Leibniz rule.  Dually, the internal Hom complex $\underline{\mathrm{Hom}}(C,D)$ is given by
[
\underline{\mathrm{Hom}}(C,D)*n = \prod_m \mathrm{Hom}(C_m,D*{m+n}),
]
with the standard differential. This is the backbone of enriched homotopical structure.

\section{Exact Sequences and Derived Functors}

A short exact sequence of complexes
[
0\to A_\bullet\to B_\bullet\to C_\bullet\to0
]
induces a long exact homology sequence. Derived functors such as $\operatorname{Ext}$ and $\operatorname{Tor}$ arise by taking projective or injective resolutions of complexes.

\begin{remark}
For a derived stack $\mathcal X$, its cotangent complex $\mathbb{L}*\mathcal X$ and tangent complex $\mathbb{T}*\mathcal X$ are defined using homological algebra in suitable model categories.
\end{remark}

\section{Model Categories}

A model category is a category equipped with three distinguished classes of morphisms (weak equivalences, fibrations, cofibrations) satisfying appropriate axioms. Homotopy categories and derived functors are constructed by formally inverting weak equivalences.

\begin{definition}
A Quillen adjunction between model categories induces an adjunction on homotopy categories, and a Quillen equivalence induces an equivalence of homotopy categories.
\end{definition}

Model category language underlies much of the theory of derived algebraic geometry and is assumed in the main text.  The reader seeking more detail will find standard references in the bibliography.

\section{Spectral Sequences}

Spectral sequences are computational tools associated to filtered chain complexes. They provide successive approximations to homology and cohomology and are indispensable for computing deformation and obstruction theories in derived geometry.

\begin{remark}
The obstruction theory used to construct the moduli stack $\Plenum$ depends on the vanishing of certain classes in higher derived $\operatorname{Ext}$ spaces, often computed via spectral sequence arguments.
\end{remark}

\section{End of Appendix B}

We emphasize that the natural home for RSVP is the $(\infty,1)$--categorical world introduced in Chapter~2, in which many of the above constructions are given a homotopical rather than merely algebraic interpretation.

